\section{2025 Analysis \& Differential Equations}

\subsection{Exam}

\begin{exercise}
\label{analysis:2025:1}
Find the solution $R(r)$ ($r\geq 1$) of the differential equation
\[
r^3 R'''(r) + 2r^2 R''(r) - rR'(r) + R(r) = 2025,\quad r>1,
\]
satisfying the conditions
\[
R(1) = 2025,\quad R'(1) = 0 \quad \text{and}\quad R''(1) = 1.
\]
\end{exercise}

\begin{solution}
\textbf{Prerequisites:}
\begin{itemize}
    \item \textbf{Euler-Cauchy Equations:} Linear differential equations where the degree of the monomial $r^k$ matches the order of the derivative $R^{(k)}$.
    \item \textbf{Change of Variables:} The substitution $r = e^t$ (or $t = \ln r$) converts Euler-Cauchy equations into linear ODEs with constant coefficients.
    \item \textbf{Chain Rule:} Used to transform derivatives from the $r$-domain to the $t$-domain.
    \item \textbf{Linear ODE Theory:} Methods for solving homogeneous equations (characteristic roots) and non-homogeneous equations (undetermined coefficients).
\end{itemize}

\bigskip

\textbf{Intuition:}
The given equation is a classic Euler-Cauchy equation. Notice how the power of $r$ exactly matches the "order" of the derivative it's attached to (e.g., $r^3$ with $R'''$). This tells us that the equation is "scale-invariant." By switching to a logarithmic scale ($t = \ln r$), we can flatten this scaling, turning it into a much simpler constant-coefficient equation that we already know how to solve using standard algebraic methods.

\bigskip

\textbf{Step 1: Transform the independent variable.}
\par
Let $r = e^t$, which means $t = \ln r$. We define a new function $y(t) = R(e^t)$. To rewrite the original ODE, we must express the derivatives of $R$ with respect to $r$ in terms of derivatives of $y$ with respect to $t$.

\par\bigskip
Using the chain rule repeatedly:
\begin{align*}
R'(r) &= \frac{dy}{dt} \frac{dt}{dr} = y'(t) \cdot \frac{1}{r} \implies rR'(r) = y'(t) \\
R''(r) &= \frac{d}{dr} \left( \frac{1}{r} y'(t) \right) = -\frac{1}{r^2} y'(t) + \frac{1}{r} \frac{d}{dr}[y'(t)] = -\frac{1}{r^2} y'(t) + \frac{1}{r^2} y''(t) = \frac{1}{r^2} (y''(t) - y'(t)) \\
R'''(r) &= \frac{d}{dr} \left( \frac{1}{r^2} (y'' - y') \right) = -\frac{2}{r^3}(y'' - y') + \frac{1}{r^3}(y''' - y'') = \frac{1}{r^3}(y''' - 3y'' + 2y')
\end{align*}

\bigskip

\textbf{Step 2: Substitute into the differential equation.}
\par
Substitute the expressions for $r^3 R'''$, $r^2 R''$, and $r R'$ into the original equation $r^3 R''' + 2r^2 R'' - rR' + R = 2025$:
\[
(y''' - 3y'' + 2y') + 2(y'' - y') - (y') + y = 2025
\]
Simplifying the coefficients for each derivative:
\[
y''' + (-3 + 2)y'' + (2 - 2 - 1)y' + y = 2025 \implies y''' - y'' - y' + y = 2025
\]

\bigskip

\textbf{Step 3: Solve the homogeneous equation.}
\par
First, we solve the homogeneous part $y''' - y'' - y' + y = 0$. Assuming a solution $y = e^{mt}$, we get:
\[
m^3 - m^2 - m + 1 = 0 \implies m^2(m - 1) - (m - 1) = 0 \implies (m^2 - 1)(m - 1) = 0 \implies (m - 1)^2(m + 1) = 0
\]
The roots are $m = 1$ (double root) and $m = -1$ (single root).
The general homogeneous solution is:
\[
y_h(t) = c_1 e^{-t} + c_2 e^t + c_3 t e^t
\]

\bigskip

\textbf{Step 4: Find a particular solution.}
\par
For the non-homogeneous equation with a constant RHS of $2025$, we guess a constant solution $y_p(t) = A$.
Plugging this in: $0 - 0 - 0 + A = 2025 \implies A = 2025$.
The total solution for $y(t)$ is:
\[
y(t) = c_1 e^{-t} + c_2 e^t + c_3 t e^t + 2025
\]

\bigskip

\textbf{Step 5: Convert back to $r$ and apply initial conditions.}
\par
Replace $e^t = r$ and $t = \ln r$:
\[
R(r) = \frac{c_1}{r} + c_2 r + c_3 r \ln r + 2025
\]
Now we find the constants using $R(1) = 2025, R'(1) = 0, R''(1) = 1$:
\begin{align*}
R'(r) &= -\frac{c_1}{r^2} + c_2 + c_3(\ln r + 1) \\
R''(r) &= \frac{2c_1}{r^3} + \frac{c_3}{r}
\end{align*}
Setting $r = 1$:
1. $R(1) = c_1 + c_2 + 2025 = 2025 \implies c_1 + c_2 = 0$.
2. $R'(1) = -c_1 + c_2 + c_3 = 0$.
3. $R''(1) = 2c_1 + c_3 = 1$.
From (1), $c_2 = -c_1$. Substituting into (2): $-2c_1 + c_3 = 0 \implies c_3 = 2c_1$.
Substituting into (3): $2c_1 + 2c_1 = 1 \implies 4c_1 = 1 \implies c_1 = 1/4$.
Then $c_2 = -1/4$ and $c_3 = 1/2$.

\par\bigskip
\textbf{Final Result:}
Substituting the constants back into $R(r)$:
\[
R(r) = \frac{1}{4r} - \frac{r}{4} + \frac{r \ln r}{2} + 2025
\]
\end{solution}

\begin{exercise}
\label{analysis:2025:2}
\begin{itemize}
\item[(1)] Prove that for $n\geq 3$, there exists a constant $C>0$ such that
\[
\int_{\mathbb{R}^n}\frac{u^2}{|x|^2}\,dx \leq C\int_{\mathbb{R}^n}|\nabla u|^2\,dx, \quad \forall u\in H^1(\mathbb{R}^n);
\]
\item[(2)] Prove that the optimal constant $C$ in (1) is $\frac{4}{(n-2)^2}$.
\end{itemize}
\end{exercise}

\begin{solution}
\textbf{Motivation:} This inequality is known as Hardy's inequality. It is extremely useful in partial differential equations because it allows us to control the $L^2$ norm of a function $u$ (weighted heavily near the origin by $1/|x|^2$) using only the $L^2$ norm of its gradient $\nabla u$.

\textbf{Part (1): Proving the inequality}
First, it is a standard result in Sobolev spaces that smooth functions with compact support, denoted $C_c^\infty(\mathbb{R}^n)$, are dense in the Sobolev space $H^1(\mathbb{R}^n)$. Because the origin $x=0$ is just a single point (and $n \geq 3$), we can further assume our test functions $u$ vanish near the origin, meaning $u \in C_c^\infty(\mathbb{R}^n \setminus \{0\})$. If we prove the inequality for these smooth functions, it holds for all $H^1$ functions by taking limits.

The trick to proving Hardy's inequality is to use integration by parts with a cleverly chosen vector field. Let's define the vector field $\mathbf{F}(x) = \frac{x}{|x|^2}$.
We first compute the divergence of $\mathbf{F}$, denoted as $\text{div} \mathbf{F}$ or $\nabla \cdot \mathbf{F}$.
Recall that $x = (x_1, x_2, \dots, x_n)$ and $|x|^2 = x_1^2 + x_2^2 + \dots + x_n^2$.
The $i$-th component of $\mathbf{F}$ is $F_i = \frac{x_i}{|x|^2}$. The partial derivative is:
\[
\partial_i F_i = \frac{\partial}{\partial x_i} \left( x_i (|x|^2)^{-1} \right) = 1 \cdot (|x|^2)^{-1} + x_i \cdot (-1)(|x|^2)^{-2} \cdot 2x_i = \frac{1}{|x|^2} - \frac{2x_i^2}{|x|^4}.
\]
Summing over all $i$ from 1 to $n$:
\[
\text{div} \mathbf{F} = \sum_{i=1}^n \left( \frac{1}{|x|^2} - \frac{2x_i^2}{|x|^4} \right) = \frac{n}{|x|^2} - \frac{2 \sum x_i^2}{|x|^4} = \frac{n}{|x|^2} - \frac{2|x|^2}{|x|^4} = \frac{n-2}{|x|^2}.
\]

Now, we evaluate the integral of $\frac{u^2}{|x|^2}$. Using our divergence identity, we have:
\[
\int_{\mathbb{R}^n} \frac{u^2}{|x|^2} dx = \frac{1}{n-2} \int_{\mathbb{R}^n} u^2 \text{div} \mathbf{F} dx.
\]
We apply integration by parts (the Divergence Theorem). Since $u$ has compact support, there are no boundary terms at infinity:
\[
\int_{\mathbb{R}^n} u^2 \text{div} \mathbf{F} dx = - \int_{\mathbb{R}^n} \nabla(u^2) \cdot \mathbf{F} dx = - \int_{\mathbb{R}^n} (2u \nabla u) \cdot \frac{x}{|x|^2} dx.
\]
Taking the absolute value of both sides and using the property that $|x \cdot y| \leq |x||y|$:
\[
(n-2) \int_{\mathbb{R}^n} \frac{u^2}{|x|^2} dx \leq 2 \int_{\mathbb{R}^n} |u| |\nabla u| \frac{|x|}{|x|^2} dx = 2 \int_{\mathbb{R}^n} \left( \frac{|u|}{|x|} \right) |\nabla u| dx.
\]
Now we apply the Cauchy-Schwarz inequality for integrals, $\int |fg| \leq \left(\int f^2\right)^{1/2} \left(\int g^2\right)^{1/2}$, with $f = \frac{|u|}{|x|}$ and $g = |\nabla u|$:
\[
(n-2) \int_{\mathbb{R}^n} \frac{u^2}{|x|^2} dx \leq 2 \left( \int_{\mathbb{R}^n} \frac{u^2}{|x|^2} dx \right)^{1/2} \left( \int_{\mathbb{R}^n} |\nabla u|^2 dx \right)^{1/2}.
\]
Assume $\int \frac{u^2}{|x|^2} dx > 0$ (otherwise the inequality is trivially $0 \leq 0$). Divide both sides by $\left( \int \frac{u^2}{|x|^2} dx \right)^{1/2}$:
\[
(n-2) \left( \int_{\mathbb{R}^n} \frac{u^2}{|x|^2} dx \right)^{1/2} \leq 2 \left( \int_{\mathbb{R}^n} |\nabla u|^2 dx \right)^{1/2}.
\]
Squaring both sides and isolating the integral of $u^2/|x|^2$:
\[
(n-2)^2 \int_{\mathbb{R}^n} \frac{u^2}{|x|^2} dx \leq 4 \int_{\mathbb{R}^n} |\nabla u|^2 dx \implies \int_{\mathbb{R}^n} \frac{u^2}{|x|^2} dx \leq \frac{4}{(n-2)^2} \int_{\mathbb{R}^n} |\nabla u|^2 dx.
\]
This proves the inequality with $C = \frac{4}{(n-2)^2}$.

\textbf{Part (2): Proving the constant is optimal}
To show $C = \frac{4}{(n-2)^2}$ is the best possible constant, we must show that no smaller constant works. We do this by testing the inequality with a specific sequence of functions.
Let's consider a radial function $u(x) = |x|^{-\alpha}$. The gradient is $\nabla u = -\alpha |x|^{-\alpha - 1} \frac{x}{|x|}$, so $|\nabla u|^2 = \alpha^2 |x|^{-2\alpha - 2}$.
If we plug this into Hardy's inequality:
\[
\int_{\mathbb{R}^n} \frac{(|x|^{-\alpha})^2}{|x|^2} dx \leq C \int_{\mathbb{R}^n} \alpha^2 |x|^{-2\alpha - 2} dx \implies \int_{\mathbb{R}^n} |x|^{-2\alpha - 2} dx \leq C \alpha^2 \int_{\mathbb{R}^n} |x|^{-2\alpha - 2} dx.
\]
Assuming the integral is non-zero and finite, we can divide both sides by it, leaving:
\[
1 \leq C \alpha^2 \implies C \geq \frac{1}{\alpha^2}.
\]
To make $C$ as small as possible, we need to make $\alpha$ as large as possible. However, we cannot choose any $\alpha$. For $u(x)$ to be in $L^2_{loc}$ (square-integrable near the origin), we need the integral $\int_0^R r^{-2\alpha} r^{n-1} dr$ to converge, which requires $-2\alpha + n - 1 > -1$, meaning $\alpha < n/2$.
Similarly, for the gradient to be in $L^2$, we need $\int_0^R r^{-2\alpha-2} r^{n-1} dr$ to converge, which requires $-2\alpha - 2 + n - 1 > -1 \implies \alpha < (n-2)/2$.
The limiting case where the function *just barely* fails to be in $H^1$ is when $\alpha = \frac{n-2}{2}$.
By defining a sequence of smooth cutoff functions that approximate $u(x) = |x|^{-(n-2)/2}$, the ratio of the two integrals will approach $\frac{1}{\alpha^2} = \frac{1}{((n-2)/2)^2} = \frac{4}{(n-2)^2}$.
Thus, no constant smaller than $\frac{4}{(n-2)^2}$ can satisfy the inequality for all $u \in H^1(\mathbb{R}^n)$.
\end{solution}

\begin{exercise}
\label{analysis:2025:3}
\begin{itemize}
\item[(1)] Let $u(x,y)$ be a subharmonic function on a domain $D$ containing the annulus $\{r_1<|x|<r_2\}$. Define
\[
M(r) = \max_{x^2+y^2=r} u(x,y)
\]
for $r_1<r<r_2$. Prove that:
\[
M(r) \leq \frac{M(r_1)(\log r_2 - \log r) + M(r_2)(\log r - \log r_1)}{\log r_2 - \log r_1};
\]
\item[(2)] If $u$ is subharmonic on $\mathbb{R}^2\setminus\{0\}$ and bounded above, prove that $u$ is constant.
\end{itemize}
\end{exercise}

\begin{solution}
\textbf{Motivation:} This problem is a classic application of the Maximum Principle for subharmonic functions, and the result in part (1) is known as Hadamard's Three-Circle Theorem. It shows that the maximum of a subharmonic function on concentric circles is convex with respect to the logarithm of the radius.

\textbf{Part (1): Proving Hadamard's Three-Circle Theorem}
The key idea is to compare the subharmonic function $u(x,y)$ with a carefully chosen harmonic function on the annulus.
Recall that in $\mathbb{R}^2$, the fundamental harmonic function (which depends only on the radius $r = \sqrt{x^2+y^2}$) is $\log r$. Any function of the form $L(x,y) = a \log r + b$, where $a$ and $b$ are constants, is harmonic in any region excluding the origin.

We want to choose $a$ and $b$ so that $L(x,y)$ perfectly matches $M(r_1)$ on the inner boundary (circle of radius $r_1$) and $M(r_2)$ on the outer boundary (circle of radius $r_2$).
This gives us a system of two linear equations:
\begin{align*}
a \log r_1 + b &= M(r_1) \quad \text{(Equation 1)} \\
a \log r_2 + b &= M(r_2) \quad \text{(Equation 2)}
\end{align*}
Subtracting Equation 1 from Equation 2 eliminates $b$:
\[
a(\log r_2 - \log r_1) = M(r_2) - M(r_1) \implies a = \frac{M(r_2) - M(r_1)}{\log r_2 - \log r_1}.
\]
Now, let's define a new function $v(x,y) = u(x,y) - L(x,y)$.
Because $u$ is subharmonic and $L$ is harmonic, their difference $v$ is also subharmonic.
Let's check the values of $v(x,y)$ on the boundaries of the annulus:
- On the inner circle $|(x,y)| = r_1$:
  $u(x,y) \leq M(r_1)$ by definition of $M$.
  $L(x,y) = M(r_1)$ by our construction.
  Therefore, $v(x,y) = u(x,y) - L(x,y) \leq M(r_1) - M(r_1) = 0$.
- On the outer circle $|(x,y)| = r_2$:
  Similarly, $u(x,y) \leq M(r_2)$ and $L(x,y) = M(r_2)$, so $v(x,y) \leq 0$.

By the Maximum Principle for subharmonic functions, if a subharmonic function is $\leq 0$ on the boundary of a bounded domain, it must be $\leq 0$ everywhere inside the domain. Thus, $v(x,y) \leq 0$ for all $r_1 \leq r \leq r_2$.
This implies $u(x,y) \leq L(x,y)$ for all points in the annulus. Taking the maximum of $u(x,y)$ over a circle of radius $r$ gives:
\[
M(r) \leq \max_{x^2+y^2=r} L(x,y) = a \log r + b.
\]
To get the final formula, we can substitute $a$ and $b$ back, or more simply, notice that $a \log r + b$ is the linear interpolation of the values $M(r_1)$ and $M(r_2)$ with respect to the variable $\log r$. Using the point-slope form:
\[
M(r) \leq M(r_1) + \frac{M(r_2) - M(r_1)}{\log r_2 - \log r_1} (\log r - \log r_1).
\]
Expanding and collecting terms for $M(r_1)$ and $M(r_2)$:
\[
M(r) \leq \frac{M(r_1)(\log r_2 - \log r_1) - M(r_1)(\log r - \log r_1) + M(r_2)(\log r - \log r_1)}{\log r_2 - \log r_1}.
\]
Notice that $(\log r_2 - \log r_1) - (\log r - \log r_1) = \log r_2 - \log r$. Thus:
\[
M(r) \leq \frac{M(r_1)(\log r_2 - \log r) + M(r_2)(\log r - \log r_1)}{\log r_2 - \log r_1}.
\]

\textbf{Part (2): Bounded subharmonic functions on the punctured plane}
We are given that $u$ is subharmonic on $\mathbb{R}^2 \setminus \{0\}$ and bounded above, meaning there exists a constant $K$ such that $u(x,y) \leq K$ for all $(x,y)$.
First, because $u$ is bounded above near the isolated singularity at the origin, a standard theorem states that the singularity is removable. This means we can extend $u$ to a function that is subharmonic on the entire plane $\mathbb{R}^2$ by defining $u(0) = \limsup_{(x,y) \to (0,0)} u(x,y)$.
Once $u$ is subharmonic on all of $\mathbb{R}^2$ and bounded above, Liouville's Theorem for subharmonic functions immediately tells us that $u$ must be constant.

Alternatively, we can use the result from Part (1) to prove this directly without the removability theorem.
Assume $u \leq K$. Let $r_2 \to \infty$ in the Three-Circle formula:
\[
M(r) \leq \limsup_{r_2 \to \infty} \left( \frac{M(r_1) \log(r_2/r)}{\log(r_2/r_1)} + \frac{M(r_2) \log(r/r_1)}{\log(r_2/r_1)} \right).
\]
Since $M(r_2) \leq K$, the second term is bounded by $K \frac{\log(r/r_1)}{\log(r_2/r_1)}$, which goes to 0 as $r_2 \to \infty$.
The first term coefficient $\frac{\log(r_2/r)}{\log(r_2/r_1)}$ goes to 1.
Thus, taking the limit gives $M(r) \leq M(r_1)$ for any $r > r_1$.
This means $M(r)$ is a decreasing function of $r$. However, by the Maximum Principle, the maximum of a subharmonic function on circles of radius $r$ must be a non-decreasing function of $r$.
The only way a function can be both non-decreasing and decreasing is if it is constant. Therefore, $M(r)$ is constant, and $u(x,y)$ must be constant.
\end{solution}

\begin{exercise}
\label{analysis:2025:4}
Write $\mathbb{S}^1 = \{e^{2\pi\sqrt{-1}s}\mid s\in\mathbb{R}\}\cong\{s\mid s\in[0,1]\} = [0,1]$. For a smooth closed curve
\[
\gamma:\mathbb{S}^1\longrightarrow\mathbb{R}^2,\quad s\longmapsto\gamma(s) = (x(s),y(s))
\]
we mean that $x,y:\mathbb{S}^1\to\mathbb{R}$ are continuously differentiable functions of $s$, $\gamma(0)=\gamma(1)$, and $\gamma'(s) = \frac{d\gamma(s)}{ds}\neq 0$ for all $s\in[0,1]$.

The degree of a smooth closed curve $\gamma:\mathbb{S}^1\to\mathbb{R}^2$ is given by
\[
d(\gamma) := \frac{1}{2\pi}\oint_{\gamma}\frac{x\,dy - y\,dx}{x^2+y^2}.
\]

\begin{itemize}
\item[(1)] Show that
\[
d(\gamma) = \frac{1}{2\pi\sqrt{-1}}\oint_{\gamma}\frac{dz}{z}.
\]
Hence $d(\gamma)$ is an integer.

\item[(2)] For the curve $\gamma(s) = z_0 e^{2\pi\sqrt{-1}ns}\in\mathbb{C}$ with $z_0\in\mathbb{C}$, $s\in[0,1]$ and $n\in\mathbb{Z}$, compute $d(\gamma)$.

\item[(3)] Two smooth closed curves $\gamma_0$ and $\gamma_1$ are regularly homotopic if there exists a family of smooth closed curves $\gamma_t$ with a continuous dependence on $t\in[0,1]$, where the curves $\gamma_0$ and $\gamma_1$ correspond to $t=0$ and $t=1$ respectively. Here, by a continuous dependence on $t$ we mean that the map
\[
[0,1]\times[0,1]\longrightarrow\mathbb{R}^2,\quad (s,t)\longmapsto\gamma(s,t)=\gamma_t(s)
\]
is continuous.

Prove that two smooth closed curves $\gamma_0$ and $\gamma_1$ are regularly homotopic if and only if they have equal degrees.
\end{itemize}
\end{exercise}

\begin{solution}
\textbf{Motivation:} This problem connects vector calculus (line integrals) with complex analysis (winding numbers) and algebraic topology (homotopy classes).

\textbf{Part (1): Connecting real and complex integrals}
Let's identify the Euclidean plane $\mathbb{R}^2$ with the complex plane $\mathbb{C}$ by letting $z = x + iy$.
Then the differential is $dz = dx + i dy$.
Let's evaluate the complex integral $\oint_\gamma \frac{dz}{z}$:
\[
\frac{dz}{z} = \frac{dx + i dy}{x + iy}.
\]
To separate the real and imaginary parts, we multiply the numerator and denominator by the complex conjugate of the denominator, $x - iy$:
\[
\frac{dz}{z} = \frac{(dx + i dy)(x - iy)}{(x + iy)(x - iy)} = \frac{(x dx + y dy) + i(x dy - y dx)}{x^2 + y^2}.
\]
Integrating this over the closed curve $\gamma$ gives:
\[
\oint_\gamma \frac{dz}{z} = \oint_\gamma \frac{x dx + y dy}{x^2 + y^2} + i \oint_\gamma \frac{x dy - y dx}{x^2 + y^2}.
\]
Let's look at the real part. Notice that $d(x^2 + y^2) = 2x dx + 2y dy$. So the real integral is exactly:
\[
\oint_\gamma \frac{1}{2} \frac{d(x^2 + y^2)}{x^2 + y^2} = \oint_\gamma d\left(\frac{1}{2} \ln(x^2 + y^2)\right).
\]
Because $\gamma$ is a closed curve and the origin is not on the curve (implied by the denominator), the integral of an exact differential over a closed loop is zero.
Thus, only the imaginary part remains:
\[
\oint_\gamma \frac{dz}{z} = i \oint_\gamma \frac{x dy - y dx}{x^2 + y^2}.
\]
Dividing both sides by $2\pi i$ (where $i = \sqrt{-1}$), we get:
\[
\frac{1}{2\pi \sqrt{-1}} \oint_\gamma \frac{dz}{z} = \frac{1}{2\pi} \oint_\gamma \frac{x dy - y dx}{x^2 + y^2}.
\]
By definition, the right side is $d(\gamma)$.
From complex analysis (Cauchy's Argument Principle), $\frac{1}{2\pi i} \oint_\gamma \frac{dz}{z}$ counts the number of times the curve winds around the origin counterclockwise. This winding number is always an integer.

\textbf{Part (2): Computing the degree of a specific curve}
We are given $\gamma(s) = z_0 e^{2\pi i n s}$ for $s \in [0, 1]$.
Let's parameterize the complex integral using $s$. We have $z(s) = z_0 e^{2\pi i n s}$.
Taking the derivative with respect to $s$:
\[
dz = z'(s) ds = z_0 (2\pi i n) e^{2\pi i n s} ds = 2\pi i n \cdot z(s) ds.
\]
Now substitute this into the integral formula from Part (1):
\[
d(\gamma) = \frac{1}{2\pi i} \oint_\gamma \frac{dz}{z} = \frac{1}{2\pi i} \int_0^1 \frac{2\pi i n \cdot z(s) ds}{z(s)} = \frac{1}{2\pi i} \int_0^1 2\pi i n ds.
\]
The $2\pi i$ cancels out, leaving:
\[
d(\gamma) = \int_0^1 n ds = n [s]_0^1 = n.
\]
This makes intuitive sense: the curve $e^{2\pi i n s}$ wraps around the origin $n$ times.

\textbf{Part (3): Regular homotopy and degrees}
First, suppose $\gamma_0$ and $\gamma_1$ are regularly homotopic. This means there is a continuous deformation $\gamma_t(s)$ between them such that for every $t$, $\gamma_t$ is a smooth closed curve that does not pass through the origin (its derivative is never zero, which is the definition of a regular curve). Since $d(\gamma_t)$ is an integer for all $t \in [0, 1]$, and it depends continuously on $t$ (because the integral of a smooth family of curves is continuous), the function $f(t) = d(\gamma_t)$ is a continuous function from $[0, 1]$ to $\mathbb{Z}$. The only continuous functions from an interval to the discrete set of integers are constant functions. Therefore, $d(\gamma_0) = d(\gamma_1)$.

Conversely, suppose $d(\gamma_0) = d(\gamma_1) = n$. This means both curves wind around the origin $n$ times.
In algebraic topology, maps from the circle $\mathbb{S}^1$ to the punctured plane $\mathbb{R}^2 \setminus \{0\} \simeq \mathbb{C}^*$ are completely classified by their winding number (the fundamental group $\pi_1(\mathbb{C}^*) \cong \mathbb{Z}$). Because they have the same winding number, they are homotopic in $\mathbb{C}^*$. Furthermore, because the curves are smooth immersions (their derivatives are non-zero), a stronger theorem known as the Whitney-Graustein Theorem guarantees that they are not just homotopic, but *regularly* homotopic (meaning the deformation itself consists only of smooth immersions).
\end{solution}

\begin{exercise}
\label{analysis:2025:5}
If $A$ and $B$ are disjoint closed subsets of a metric space $(X,d)$ and $[a,b]$ is a given closed interval, then there exists a continuous map $f:X\to[a,b]$ such that $f(A)=\{a\}$ and $f(B)=\{b\}$.
\end{exercise}

\begin{solution}
\textbf{Motivation:} This statement is a special case of Urysohn's Lemma. Urysohn's Lemma is a fundamental topological theorem showing that in certain spaces (like metric spaces), any two disjoint closed sets can be strictly separated by a continuous function.

\textbf{Step-by-step Proof:}
To build this function, we will use the properties of the distance metric $d$.
First, let's define the distance from a point $x \in X$ to a set $S \subset X$:
\[
d(x, S) = \inf \{ d(x, y) \mid y \in S \}.
\]
A well-known property of metric spaces is that the function $x \mapsto d(x, S)$ is always continuous. In fact, it is Lipschitz continuous.
Furthermore, because $A$ is a closed set, $d(x, A) = 0$ if and only if $x \in A$.
Similarly, because $B$ is a closed set, $d(x, B) = 0$ if and only if $x \in B$.

We want to construct a function $f(x)$ that outputs $a$ when $x \in A$ and $b$ when $x \in B$.
Let's look at the expression:
\[
\frac{d(x, A)}{d(x, A) + d(x, B)}.
\]
Let's analyze this expression carefully:
1. \textbf{Denominator is never zero:} Since $A$ and $B$ are disjoint ($A \cap B = \emptyset$), a point $x$ cannot belong to both $A$ and $B$. Thus, $d(x, A)$ and $d(x, B)$ cannot both be zero at the same time. Since distances are non-negative, the sum $d(x, A) + d(x, B)$ is strictly greater than zero for all $x \in X$.
2. \textbf{Continuity:} Because both $d(x, A)$ and $d(x, B)$ are continuous, and the denominator is never zero, the entire fraction is a continuous function.
3. \textbf{Behavior on $A$:} If $x \in A$, then $d(x, A) = 0$. The fraction becomes $\frac{0}{0 + d(x, B)} = 0$.
4. \textbf{Behavior on $B$:} If $x \in B$, then $d(x, B) = 0$. The fraction becomes $\frac{d(x, A)}{d(x, A) + 0} = 1$.
5. \textbf{Range:} Since $d(x, A)$ and $d(x, B)$ are non-negative, the numerator is always less than or equal to the denominator, so the fraction always outputs a value in the interval $[0, 1]$.

Now, we just need to scale and shift this function so that its range is $[a, b]$ instead of $[0, 1]$.
We define our final continuous function $f: X \to [a, b]$ as:
\[
f(x) = a + (b - a) \frac{d(x, A)}{d(x, A) + d(x, B)}.
\]
Let's verify it meets all conditions:
- \textbf{Continuity:} It is a linear transformation of a continuous function, so it is continuous.
- \textbf{Value on $A$:} For $x \in A$, the fraction is 0, so $f(x) = a + (b-a)(0) = a$.
- \textbf{Value on $B$:} For $x \in B$, the fraction is 1, so $f(x) = a + (b-a)(1) = b$.
- \textbf{Range:} Since the fraction is between 0 and 1, $f(x)$ ranges exactly from $a$ to $a + (b-a) = b$.

This completely proves the existence of such a function.
\end{solution}

\begin{exercise}
\label{analysis:2025:6}
For $f$ a complex-valued locally integrable function on $\mathbb{R}^n$, set
\[
\|f\|_{\mathbf{BMO}} := \sup_Q\left(\frac{1}{|Q|}\int_Q|f(x)-\mathbf{Avg}_Q f|\,dx\right),
\]
where the supremum is taken over all cubes $Q$ in $\mathbb{R}^n$ and
\[
\mathbf{Avg}_Q f := \frac{1}{|Q|}\int_Q f(x)\,dx.
\]

Let $\mathbf{BMO}(\mathbb{R}^n)$ be the set of all complex-valued locally integrable functions on $\mathbb{R}^n$ with $\|f\|_{\mathbf{BMO}}<\infty$.

\begin{enumerate}
\item[(1)] $\|\cdot\|_{\mathbf{BMO}}$ is not a norm, and is only a seminorm.
\item[(2)] $\ell^{\infty}(\mathbb{R}^n)\subset\mathbf{BMO}(\mathbb{R}^n)$ and $\|f\|_{\mathbf{BMO}} \leq 2\|f\|_{L^{\infty}}$.
\item[(3)] Suppose that there exists an $A>0$ such that for all cubes $Q$ in $\mathbb{R}^n$ there exists a constant $c_Q$ such that
\[
\sup_Q\left(\frac{1}{|Q|}\int_Q|f(x)-c_Q|\,dx\right) \leq A.
\]
Then $f\in\mathbf{BMO}(\mathbb{R}^n)$ and $\|f\|_{\mathbf{BMO}} \leq 2A$.
\item[(4)] $L^{\infty}(\mathbb{R}^n)$ is a proper subspace of $\mathbf{BMO}(\mathbb{R}^n)$.
\end{enumerate}
\end{exercise}

\begin{solution}
\textbf{Motivation:} Bounded Mean Oscillation (BMO) is a crucial space in harmonic analysis. It contains functions that might be unbounded (unlike $L^\infty$), but their average deviation from their local mean is bounded.

\textbf{Part (1): Seminorm vs. Norm}
A norm $\|\cdot\|$ must satisfy the property that $\|f\| = 0$ if and only if $f$ is the zero function ($f(x) = 0$ almost everywhere).
Let's test this property on the BMO definition. Consider a non-zero constant function, $f(x) = C \neq 0$ for all $x$.
For any cube $Q$, the average of $f$ is:
\[
\mathbf{Avg}_Q f = \frac{1}{|Q|} \int_Q C \,dx = C.
\]
Then, the deviation $|f(x) - \mathbf{Avg}_Q f|$ is exactly $|C - C| = 0$ everywhere.
Therefore, the integral of the deviation is 0, and taking the supremum over all cubes yields $\|f\|_{\mathbf{BMO}} = 0$.
Because a non-zero function has a BMO "norm" of 0, $\|\cdot\|_{\mathbf{BMO}}$ fails the strict positivity requirement. It is therefore only a seminorm. (To make it a true norm, we would have to consider the space of functions modulo constants).

\textbf{Part (2): BMO contains $L^\infty$}
Suppose $f \in L^\infty(\mathbb{R}^n)$, which means $|f(x)| \leq \|f\|_{L^\infty}$ almost everywhere.
We want to bound the BMO expression for any arbitrary cube $Q$:
\[
|f(x) - \mathbf{Avg}_Q f|.
\]
By the triangle inequality:
\[
|f(x) - \mathbf{Avg}_Q f| \leq |f(x)| + |\mathbf{Avg}_Q f|.
\]
We know $|f(x)| \leq \|f\|_{L^\infty}$. Furthermore, the absolute value of the average is bounded by the maximum of the function:
\[
|\mathbf{Avg}_Q f| = \left| \frac{1}{|Q|} \int_Q f(y) \,dy \right| \leq \frac{1}{|Q|} \int_Q |f(y)| \,dy \leq \frac{1}{|Q|} \int_Q \|f\|_{L^\infty} \,dy = \|f\|_{L^\infty}.
\]
Thus, the integrand is bounded by:
\[
|f(x) - \mathbf{Avg}_Q f| \leq \|f\|_{L^\infty} + \|f\|_{L^\infty} = 2\|f\|_{L^\infty}.
\]
Now, plug this bound back into the BMO integral for the cube $Q$:
\[
\frac{1}{|Q|} \int_Q |f(x) - \mathbf{Avg}_Q f| \,dx \leq \frac{1}{|Q|} \int_Q 2\|f\|_{L^\infty} \,dx = 2\|f\|_{L^\infty}.
\]
Since this holds for *every* cube $Q$, the supremum over all cubes is also bounded by $2\|f\|_{L^\infty}$.
Thus, $\|f\|_{\mathbf{BMO}} \leq 2\|f\|_{L^\infty}$, which proves that every $L^\infty$ function is also a BMO function.

\textbf{Part (3): Any constant $c_Q$ works}
We are given that for any cube $Q$, there is some constant $c_Q$ such that the average deviation from $c_Q$ is $\leq A$.
We need to bound the deviation from the actual average, $\mathbf{Avg}_Q f$.
Using the trick of adding and subtracting $c_Q$ inside the absolute value, and applying the triangle inequality:
\[
|f(x) - \mathbf{Avg}_Q f| = |(f(x) - c_Q) + (c_Q - \mathbf{Avg}_Q f)| \leq |f(x) - c_Q| + |c_Q - \mathbf{Avg}_Q f|.
\]
Let's analyze the second term. By definition of the average:
\[
|c_Q - \mathbf{Avg}_Q f| = \left| c_Q - \frac{1}{|Q|} \int_Q f(y) \,dy \right| = \left| \frac{1}{|Q|} \int_Q (c_Q - f(y)) \,dy \right|.
\]
Moving the absolute value inside the integral (which can only increase the value):
\[
|c_Q - \mathbf{Avg}_Q f| \leq \frac{1}{|Q|} \int_Q |c_Q - f(y)| \,dy.
\]
Notice that this right-hand side is exactly the integral given in the problem statement, which is $\leq A$.
So, $|c_Q - \mathbf{Avg}_Q f| \leq A$.
Now we integrate our initial triangle inequality over the cube $Q$:
\begin{align*}
\frac{1}{|Q|} \int_Q |f(x) - \mathbf{Avg}_Q f| \,dx &\leq \frac{1}{|Q|} \int_Q \left( |f(x) - c_Q| + |c_Q - \mathbf{Avg}_Q f| \right) \,dx \\
&\leq \left( \frac{1}{|Q|} \int_Q |f(x) - c_Q| \,dx \right) + \left( \frac{1}{|Q|} \int_Q A \,dx \right).
\end{align*}
The first term is $\leq A$ by hypothesis. The second term integrates the constant $A$, yielding exactly $A$.
Thus, the sum is $\leq A + A = 2A$.
Taking the supremum over all cubes gives $\|f\|_{\mathbf{BMO}} \leq 2A$.

\textbf{Part (4): $L^\infty$ is a proper subspace}
To show it is a *proper* subspace, we must find a function that is in BMO but NOT in $L^\infty$.
Consider the function $f(x) = \ln |x|$ on $\mathbb{R}^n$.
- Is it in $L^\infty$? No, because as $x \to 0$, $\ln |x| \to -\infty$, and as $x \to \infty$, $\ln |x| \to \infty$. It is unbounded.
- Is it in BMO? Yes. Proving this rigorously takes some work, but intuitively, $\ln |x|$ only grows logarithmically. On any given cube, the function doesn't oscillate wildly; its average deviation from its mean remains bounded by a constant regardless of how close the cube is to the origin or infinity (this scales nicely because $\ln |cx| = \ln c + \ln |x|$).
Because such a function exists, $L^\infty(\mathbb{R}^n) \subsetneq \mathbf{BMO}(\mathbb{R}^n)$.
\end{solution}